\IfFileExists{stacks-project.cls}{%
\documentclass{stacks-project}
}{%
\documentclass{amsart}
}
\usepackage{amssymb}

% For dealing with references we use the comment environment
\usepackage{verbatim}
\newenvironment{reference}{\comment}{\endcomment}
%\newenvironment{reference}{}{}
\newenvironment{slogan}{\comment}{\endcomment}
\newenvironment{history}{\comment}{\endcomment}

% For commutative diagrams we use Xy-pic
\usepackage[all]{xy}

% We use 2cell for 2-commutative diagrams.
\xyoption{2cell}
\UseAllTwocells

% We use multicol for the list of chapters between chapters
\usepackage{multicol}

% This is generally recommended for better output
\usepackage{lmodern}
\usepackage[T1]{fontenc}
\usepackage{textalpha}

% For cross-file-references
\usepackage{xr-hyper}
\newcommand{\maybeexternaldocument}[2]{%
  \IfFileExists{#2.aux}{%
    \externaldocument[#1]{#2}%
  }{%
    \PackageWarning{stack}{Missing #2.aux; external references with prefix #1 unavailable}%
  }%
}
\let\externaldocumentorig\externaldocument
\renewcommand{\externaldocument}[2][]{%
  \IfFileExists{#2.aux}{%
    \externaldocumentorig[#1]{#2}%
  }{%
    \PackageWarning{stack}{Missing #2.aux; external references with prefix #1 unavailable}%
  }%
}

% Package for hypertext links:
\usepackage{hyperref}

% For any local file, say "hello.tex" you want to link to please
% use \externaldocument[hello-]{hello}
\externaldocument[introduction-]{introduction}
\externaldocument[conventions-]{conventions}
\externaldocument[sets-]{sets}
\externaldocument[categories-]{categories}
\externaldocument[topology-]{topology}
\externaldocument[sheaves-]{sheaves}
\externaldocument[sites-]{sites}
\externaldocument[stacks-]{stacks}
\externaldocument[fields-]{fields}
\externaldocument[algebra-]{algebra}
\externaldocument[brauer-]{brauer}
\externaldocument[homology-]{homology}
\externaldocument[derived-]{derived}
\externaldocument[simplicial-]{simplicial}
\externaldocument[more-algebra-]{more-algebra}
\externaldocument[smoothing-]{smoothing}
\externaldocument[modules-]{modules}
\externaldocument[sites-modules-]{sites-modules}
\externaldocument[injectives-]{injectives}
\externaldocument[cohomology-]{cohomology}
\externaldocument[sites-cohomology-]{sites-cohomology}
\externaldocument[dga-]{dga}
\externaldocument[dpa-]{dpa}
\externaldocument[sdga-]{sdga}
\externaldocument[hypercovering-]{hypercovering}
\externaldocument[schemes-]{schemes}
\externaldocument[constructions-]{constructions}
\externaldocument[properties-]{properties}
\externaldocument[morphisms-]{morphisms}
\externaldocument[coherent-]{coherent}
\externaldocument[divisors-]{divisors}
\externaldocument[limits-]{limits}
\externaldocument[varieties-]{varieties}
\externaldocument[topologies-]{topologies}
\externaldocument[descent-]{descent}
\externaldocument[perfect-]{perfect}
\externaldocument[more-morphisms-]{more-morphisms}
\externaldocument[flat-]{flat}
\externaldocument[groupoids-]{groupoids}
\externaldocument[more-groupoids-]{more-groupoids}
\externaldocument[etale-]{etale}
\externaldocument[chow-]{chow}
\externaldocument[intersection-]{intersection}
\externaldocument[pic-]{pic}
\externaldocument[weil-]{weil}
\externaldocument[adequate-]{adequate}
\externaldocument[dualizing-]{dualizing}
\externaldocument[duality-]{duality}
\externaldocument[discriminant-]{discriminant}
\externaldocument[derham-]{derham}
\externaldocument[local-cohomology-]{local-cohomology}
\externaldocument[algebraization-]{algebraization}
\externaldocument[curves-]{curves}
\externaldocument[resolve-]{resolve}
\externaldocument[models-]{models}
\externaldocument[functors-]{functors}
\externaldocument[equiv-]{equiv}
\externaldocument[pione-]{pione}
\externaldocument[etale-cohomology-]{etale-cohomology}
\externaldocument[proetale-]{proetale}
\externaldocument[relative-cycles-]{relative-cycles}
\externaldocument[more-etale-]{more-etale}
\externaldocument[trace-]{trace}
\externaldocument[crystalline-]{crystalline}
\externaldocument[spaces-]{spaces}
\externaldocument[spaces-properties-]{spaces-properties}
\externaldocument[spaces-morphisms-]{spaces-morphisms}
\externaldocument[decent-spaces-]{decent-spaces}
\externaldocument[spaces-cohomology-]{spaces-cohomology}
\externaldocument[spaces-limits-]{spaces-limits}
\externaldocument[spaces-divisors-]{spaces-divisors}
\externaldocument[spaces-over-fields-]{spaces-over-fields}
\externaldocument[spaces-topologies-]{spaces-topologies}
\externaldocument[spaces-descent-]{spaces-descent}
\externaldocument[spaces-perfect-]{spaces-perfect}
\externaldocument[spaces-more-morphisms-]{spaces-more-morphisms}
\externaldocument[spaces-flat-]{spaces-flat}
\externaldocument[spaces-groupoids-]{spaces-groupoids}
\externaldocument[spaces-more-groupoids-]{spaces-more-groupoids}
\externaldocument[bootstrap-]{bootstrap}
\externaldocument[spaces-pushouts-]{spaces-pushouts}
\externaldocument[spaces-chow-]{spaces-chow}
\externaldocument[groupoids-quotients-]{groupoids-quotients}
\externaldocument[spaces-more-cohomology-]{spaces-more-cohomology}
\externaldocument[spaces-simplicial-]{spaces-simplicial}
\externaldocument[spaces-duality-]{spaces-duality}
\externaldocument[formal-spaces-]{formal-spaces}
\externaldocument[restricted-]{restricted}
\externaldocument[spaces-resolve-]{spaces-resolve}
\externaldocument[formal-defos-]{formal-defos}
\externaldocument[defos-]{defos}
\externaldocument[cotangent-]{cotangent}
\externaldocument[examples-defos-]{examples-defos}
\externaldocument[algebraic-]{algebraic}
\externaldocument[examples-stacks-]{examples-stacks}
\externaldocument[stacks-sheaves-]{stacks-sheaves}
\externaldocument[criteria-]{criteria}
\externaldocument[artin-]{artin}
\externaldocument[quot-]{quot}
\externaldocument[stacks-properties-]{stacks-properties}
\externaldocument[stacks-morphisms-]{stacks-morphisms}
\externaldocument[stacks-limits-]{stacks-limits}
\externaldocument[stacks-cohomology-]{stacks-cohomology}
\externaldocument[stacks-perfect-]{stacks-perfect}
\externaldocument[stacks-introduction-]{stacks-introduction}
\externaldocument[stacks-more-morphisms-]{stacks-more-morphisms}
\externaldocument[stacks-geometry-]{stacks-geometry}
\externaldocument[moduli-]{moduli}
\externaldocument[moduli-curves-]{moduli-curves}
\externaldocument[examples-]{examples}
\externaldocument[exercises-]{exercises}
\externaldocument[guide-]{guide}
\externaldocument[desirables-]{desirables}
\externaldocument[coding-]{coding}
\externaldocument[obsolete-]{obsolete}
\externaldocument[fdl-]{fdl}
\externaldocument[index-]{index}

% Heap Project documents
\externaldocument[linear-algebra-]{linear-algebra}
\externaldocument[groups-]{groups}
\externaldocument[complex-analysis-]{complex-analysis}
\externaldocument[functional-analysis-]{functional-analysis}
\externaldocument[categories-]{categories}
\externaldocument[commutative-algebra-]{commutative-algebra}
\externaldocument[ringed-spaces-]{ringed-spaces}
\externaldocument[homological-algebra-]{homological-algebra}
\externaldocument[lie-algebras-]{lie-algebras}
\externaldocument[representation-theory-]{representation-theory}
\externaldocument[smooth-manifolds-]{smooth-manifolds}
\externaldocument[differential-topology-]{differential-topology}
\externaldocument[algebraic-topology-]{algebraic-topology}
\externaldocument[differential-geometry-]{differential-geometry}
\externaldocument[algebraic-geometry-]{algebraic-geometry}
\externaldocument[symplectic-geometry-]{symplectic-geometry}
\externaldocument[complex-geometry-]{complex-geometry}
\externaldocument[hodge-theory-]{hodge-theory}
\maybeexternaldocument{deformations-}{deformations}
\externaldocument[vertex-algebras-]{vertex-algebras}
\externaldocument[springer-]{springer}
\externaldocument[infty-categories-]{infty-categories}
\externaldocument[seminars-differential-geometry-]{%
seminars-differential-geometry}
\externaldocument[seminars-complex-geometry-]{%
seminars-complex-geometry}
\externaldocument[seminars-algebraic-geometry-]{%
seminars-algebraic-geometry}
\externaldocument[seminars-others-]{%
seminars-others}
%\externaldocument[geometric-prequantization-]{geometric-prequantization}
\externaldocument[surfaces-]{surfaces}
\externaldocument[birational-maps-commutative-algebra-]{birational-maps-commutative-algebra}
\externaldocument[cremona-transformations-]{cremona-transformations}
\externaldocument[derived-categories-of-sheaves-]{derived-categories-of-sheaves}
\externaldocument[geometric-invariant-theory-]{geometric-invariant-theory}
\externaldocument[geometric-stability-conditions-and-group-actions-]{geometric-stability-conditions-and-group-actions}
\externaldocument[moduli-spaces-of-sheaves-]{moduli-spaces-of-sheaves}
\externaldocument[bridgeland-stability-general-theory-]{bridgeland-stability-general-theory}
\externaldocument[hilbert-schemes-]{hilbert-schemes}
\externaldocument[noncommutative-abelian-surfaces-and-kummer-type-hyperkahler-varieties-]{noncommutative-abelian-surfaces-and-kummer-type-hyperkahler-varieties}
\externaldocument[mumford-tate-groups-in-hodge-theory-]{mumford-tate-groups-in-hodge-theory}
\externaldocument[hodge-birational-atoms-2026-]{hodge-birational-atoms-2026}
\externaldocument[xxii-egd-2026-]{%
xxii-egd-2026}
\externaldocument[pde-]{pde}
\externaldocument[probability-]{probability}

\externaldocument[linguistics-]{linguistics}
\externaldocument[genealogia-familiar-]{%
genealogia-familiar}

\externaldocument[%
16th-alga-meeting-2026-]{%
16th-alga-meeting-2026}

\externaldocument[%
iv-jornadas-lefschetz-2026-]{%
iv-jornadas-lefschetz-2026}

\externaldocument[reading-the-masters-]{%
reading-the-masters}

\externaldocument[real-analysis-]{%
real-analysis}

\externaldocument[optimal-transport-]{%
optimal-transport}

\externaldocument[history-]{history}

% Theorem environments.
%
\theoremstyle{plain}
\newtheorem{theorem}[subsection]{Theorem}
\newtheorem{proposition}[subsection]{Proposition}
\newtheorem{lemma}[subsection]{Lemma}

\theoremstyle{definition}
\newtheorem{definition}[subsection]{Definition}
\newtheorem{example}[subsection]{Example}
\newtheorem{exercise}[subsection]{Exercise}
\newtheorem{situation}[subsection]{Situation}

\theoremstyle{remark}
\newtheorem{remark}[subsection]{Remark}
\newtheorem{remarks}[subsection]{Remarks}

\numberwithin{equation}{subsection}

% Macros
%
\def\lim{\mathop{\mathrm{lim}}\nolimits}
\def\colim{\mathop{\mathrm{colim}}\nolimits}
\def\Spec{\mathop{\mathrm{Spec}}}
\def\Hom{\mathop{\mathrm{Hom}}\nolimits}
\def\Ext{\mathop{\mathrm{Ext}}\nolimits}
\def\SheafHom{\mathop{\mathcal{H}\!\mathit{om}}\nolimits}
\def\SheafExt{\mathop{\mathcal{E}\!\mathit{xt}}\nolimits}
\def\Sch{\mathit{Sch}}
\def\Mor{\mathop{\mathrm{Mor}}\nolimits}
\def\Ob{\mathop{\mathrm{Ob}}\nolimits}
\def\Sh{\mathop{\mathit{Sh}}\nolimits}
\def\NL{\mathop{N\!L}\nolimits}
\def\CH{\mathop{\mathrm{CH}}\nolimits}
\def\proetale{{pro\text{-}\acute{e}tale}}
\def\etale{{\acute{e}tale}}
\def\QCoh{\mathit{QCoh}}
\def\Ker{\mathop{\mathrm{Ker}}}
\def\Im{\mathop{\mathrm{Im}}}
\def\Coker{\mathop{\mathrm{Coker}}}
\def\Coim{\mathop{\mathrm{Coim}}}

% Boxtimes
%
\DeclareMathSymbol{\boxtimes}{\mathbin}{AMSa}{"02}

%
% Macros for moduli stacks/spaces
%
\def\QCohstack{\mathcal{QC}\!\mathit{oh}}
\def\Cohstack{\mathcal{C}\!\mathit{oh}}
\def\Spacesstack{\mathcal{S}\!\mathit{paces}}
\def\Quotfunctor{\mathrm{Quot}}
\def\Hilbfunctor{\mathrm{Hilb}}
\def\Curvesstack{\mathcal{C}\!\mathit{urves}}
\def\Polarizedstack{\mathcal{P}\!\mathit{olarized}}
\def\Complexesstack{\mathcal{C}\!\mathit{omplexes}}
% \Pic is the operator that assigns to X its picard group, usage \Pic(X)
% \Picardstack_{X/B} denotes the Picard stack of X over B
% \Picardfunctor_{X/B} denotes the Picard functor of X over B
\def\Pic{\mathop{\mathrm{Pic}}\nolimits}
\def\Picardstack{\mathcal{P}\!\mathit{ic}}
\def\Picardfunctor{\mathrm{Pic}}
\def\Deformationcategory{\mathcal{D}\!\mathit{ef}}


\begin{document}

\title{Arte moderna brasileira}
\maketitle

\phantomsection
\label{section-phantom}

\tableofcontents

\section{Abstração e arte brasileira no pós-guerra}

\subsection{Um mapa não linear}

A história da arte brasileira do pós-guerra não deve ser entendida como uma
sequência linear em que figuração, abstração, concretismo e neoconcretismo se
substituem. Em particular, ``figuração'' não é o nome de um primeiro movimento:
é uma categoria ampla para obras que representam figuras, objetos ou espaços
reconhecíveis.

Nos anos 1940--1960 coexistem diferentes vertentes abstratas. Duas famílias
úteis para organizar o período são a abstração informal e a abstração
geométrica. A primeira privilegia gesto, matéria, ritmo, cor e subjetividade; a
segunda privilegia relações formais e construção geométrica. Dentro desta
última se desenvolve o concretismo brasileiro nos anos 1950 e, em diálogo e
ruptura com ele, o neoconcretismo a partir de 1959.

\subsection{Concretismo e neoconcretismo}

A arte concreta busca uma linguagem visual autônoma, fundada na ordem, na
racionalidade e na construção objetiva da forma. No Brasil, um núcleo decisivo
foi o Grupo Ruptura, em São Paulo. No Rio de Janeiro, o Grupo Frente teve Ivan
Serpa como figura central e uma prática mais heterodoxa dentro do campo
geométrico.

O neoconcretismo reage ao caráter excessivamente rígido e racionalista do
programa concreto. Conserva a investigação geométrica, mas introduz uma relação
mais sensível, fenomenológica, espacial e participativa com a obra. Lygia
Clark, Lygia Pape e Hélio Oiticica são nomes centrais desta passagem. Uma
fórmula útil é: no concretismo, geometria construída; no neoconcretismo,
geometria experimentada.

Ivan Serpa é particularmente importante como figura de ligação. Além de sua
produção, sua atividade docente no MAM Rio ajudou a formar o ambiente artístico
carioca do qual emergiram experiências neoconcretas e posteriores.

Artistas como Alfredo Volpi, Maria Leontina e Milton Dacosta dialogam fortemente
com a abstração geométrica do período sem se reduzirem simplesmente aos
programas doutrinários concreto ou neoconcreto. Em Volpi, por exemplo, a
geometria das fachadas e bandeirinhas conserva uma qualidade artesanal e
pictórica muito própria.

\subsection{Abstração informal}

Contemporaneamente ao concretismo, desenvolve-se no Brasil uma abstração
informal. Aqui a tela é entendida como campo de investigação do próprio ato de
pintar: gesto, manchas, tensões, matéria e cor ganham autonomia. Essa família
tem paralelos internacionais no expressionismo abstrato norte-americano e no
Art Informel europeu.

Tomie Ohtake é uma referência importante. Sua produção dos anos 1950--1960
explora gesto, matéria e grandes relações cromáticas, em contraste com a
precisão programática da arte concreta. Algumas obras, como sua \emph{Pintura}
de 1964, podem lembrar visualmente os campos de cor de Mark Rothko, embora a
superfície de Ohtake preserve uma presença própria do gesto e da matéria.

\subsection{Genealogia internacional da abstração geométrica}

O concretismo brasileiro pertence a uma constelação internacional mais ampla de
arte abstrata e construtiva. Uma genealogia simplificada passa pelo cubismo do
início do século XX e, mais diretamente, por Malevich, Mondrian, o
construtivismo russo, De Stijl e a Bauhaus. Nos anos 1950--1960 surgem ou se
consolidam diferentes respostas locais: concretismo e neoconcretismo no Brasil,
Op Art na Europa, hard-edge painting nos Estados Unidos e tradições
construtivas importantes na Argentina e no Uruguai.

Piet Mondrian é um ancestral particularmente direto. Depois do contato com o
cubismo em Paris, sua pintura passa gradualmente da representação de árvores e
paisagens à organização autônoma de linhas, planos e cores. No neoplasticismo,
verticais e horizontais, cores primárias e relações assimétricas de equilíbrio
constituem a própria linguagem da pintura. Sua obra mostra como o cubismo é uma
das grandes rupturas que tornam possível a abstração geométrica, embora esta
não possa ser reduzida simplesmente a uma derivação do cubismo.

\subsection{Brasil e México}

Há afinidades fortes entre a abstração geométrica brasileira e a arte mexicana
do pós-guerra, mas não existe uma correspondência simples entre concretismo e
neoconcretismo brasileiros e movimentos mexicanos de mesmo tipo. No México,
a chamada Generación de la Ruptura e artistas como Vicente Rojo, Mathias
Goeritz e Gunther Gerzso participam da abertura para linguagens abstratas
internacionais, num contexto marcado também pela reação à hegemonia da Escola
Mexicana de Pintura e do muralismo.

Vicente Rojo (1932--2021), ativo no México sobretudo a partir dos anos 1950,
tem afinidades visuais com a abstração geométrica brasileira: construção em
séries, módulos, repetição, cor e uma relação importante com o design gráfico.
Maria Leontina e Milton Dacosta oferecem comparações particularmente naturais.
A relação não deve, porém, ser confundida automaticamente com contato pessoal
entre esses artistas.

Um vínculo histórico México--Brasil mais concreto aparece na poesia concreta.
O grupo brasileiro Noigandres participa de uma circulação internacional que
chega ao México; Mathias Goeritz organiza em 1966, na UNAM, a exposição
\emph{Poesía Concreta Internacional}, na qual Vicente Rojo participa. Assim,
geometria, tipografia, poesia, livro e design constituem uma zona real de
intercâmbio entre os dois países.

\subsection{Retorno da figuração e arte experimental}

Nos anos 1960--1980, a figuração reaparece com novas funções. A Nova Figuração
brasileira incorpora cultura de massa, imprensa, televisão e referências
cotidianas, e no contexto da ditadura militar torna-se também veículo de
crítica social e política. Não se trata de retornar à figuração anterior à
abstração: a imagem retorna depois do cubismo, da abstração e da cultura de
massa.

Ao mesmo tempo, a arte experimental e conceitual radicaliza outra pergunta. Em
vez de buscar apenas uma nova linguagem pictórica, questiona o próprio objeto
artístico, seus suportes, sua circulação e as instituições que o legitimam. A
ênfase pode deslocar-se do objeto para o conceito ou procedimento. Performance,
fotografia, intervenção urbana, corpo e ready-made passam a integrar o campo da
arte. No Brasil, esse desenvolvimento adquire uma dimensão política específica
sob a ditadura.

\subsection{Depois de 1990}

A partir dos anos 1990 torna-se menos convincente narrar a arte brasileira como
uma sucessão de movimentos estilísticos dominantes. Figuração, abstração,
pintura, práticas conceituais, performance, vídeo e meios digitais coexistem.
A revisão dos cânones, a ampliação dos agentes reconhecidos pelo sistema da arte
e questões de memória, território, corpo, identidade e pertencimento tornam-se
centrais. A ideia moderna de ruptura continua presente, mas já não produz uma
sequência tão clara de novos ``ismos''.


\section{Salvador: arte, memória e museu}
\label{section-arte-salvador-2026}

Estas notas surgiram da visita ao MUNCAB e à
Casa do Benin em 26 de setembro
de 2026. A distinção metodológica principal é
entre a obra, sua legenda, o
texto curatorial e uma interpretação
posterior. Esses níveis não devem ser
fundidos.

\section{MUNCAB e a 36a Bienal}
A itinerância da 36a Bienal de São Paulo no
MUNCAB usava o título
\emph{Nem todo viandante anda estradas---Da
humanidade como prática}. O texto
curatorial perguntava como a história da arte
produz presenças, ausências,
centralidades e silêncios, e recusava a ideia
de que pertencimento exige uma
identidade homogênea \cite{muncab2026bienal}.

Esse primeiro painel não foi o ponto mais
convincente da visita; a exposição de
Hariel Revignet tornou questões semelhantes
muito mais concretas através de
água, corpo, território e memória.

\section{Hariel Revignet: Omi Tutu}
\emph{Omi Tutu---Água fresca}, de Hariel
Revignet, organiza a água como
elemento de conexão entre territórios,
corpos, deslocamento, continuidade e
retorno. A artista é brasileira-gabonesa; o
Gabão, país da África Central na
costa atlântica, entra na exposição como
território familiar concreto e não
como uma África abstrata
\cite{muncab2026omitutu}.

A imagem curatorial mais memorável foi a do
rio que entra no mar: depois da
confluência já não se distingue o rio, mas
sua água não deixa de existir.
Isso fornece uma imagem forte para pensar
continuidade diaspórica sem supor a
sobrevivência intacta de uma origem.

As legendas \emph{Ori Ejá} e \emph{Iya Ejá},
ambas de 2025, acompanhavam
pinturas em acrílica com bordado de búzios.
Na fotografia feita durante a
visita, uma figura feminina escura aparece
com ventre arredondado, mãos sobre
o corpo e cauda de peixe. Como duas legendas
estavam adjacentes, a fotografia
sozinha não permite decidir qual título
correspondia à tela fotografada; essa
incerteza deve ser preservada.

\section{Rio de Miçangas---Ziguidás}
A instalação \emph{Rio de
Miçangas---Ziguidás} suspendia fios de
miçangas e
sementes sobre recipientes de barro, cabaças
e outros objetos, produzindo
visualmente uma espécie de rio vertical.

A relação com os zigidas/sikidas é biográfica
e material: textos sobre a
trajetória de Hariel descrevem esses fios de
contas usados abaixo do umbigo
como objetos ligados a riqueza, sensualidade
e fertilidade, transmitidos no
caso da artista através de sua avó gabonesa
\cite{hariel-zigida}. A obra
transforma assim um objeto corporal em
instalação espacial, sem que seu sentido
se reduza a uma única função ritual.

\section{Nádia Taquary e Iroko}
Nádia Taquary, em \emph{Iroko: A árvore
cósmica} (2025), usa fios de contas e
uma linguagem escultórica contemporânea para
trabalhar tempo, ancestralidade,
ciclo de vida e referências a Iroko e às
Iyami.

A obra é útil para perceber que vocabulário
religioso pode entrar na arte
contemporânea sem que a obra se torne uma
ilustração etnográfica neutra. A
tradução de Iyami por feiticeiras no painel
do museu é, ela própria, uma
decisão de tradução que deve ser analisada
criticamente.

\section{Sandro Aiyê e Bará}
A escultura \emph{Bará II}, de Sandro Aiyê,
foi identificada na legenda como
madeira de demolição e contas de miçangas.
Ela integrava
\emph{Padê Onã---Encontrar Caminhos},
exposição que o MUNCAB organiza em torno
de Exu, movimento, mediação, circulação e
abertura de caminhos
\cite{muncab2026padeona}.

Isso é mais forte do que identificar a obra
apenas por uma forma semelhante a
um tridente: título, material, exposição e
contexto curatorial trabalham
juntos.

\section{Alberto Pitta e J. Cunha}
Dois livros encontrados no MUNCAB abrem uma
linha importante para a história
da arte baiana.

\emph{Alberto Pitta: Fúnfún Dúdú} reúne a
produção de Pitta em estamparia,
pintura e objetos, articulando terreiro,
carnaval e cidade
\cite{barata2025pitta}. \emph{J. Cunha e o
Carnaval Negro} examina a
trajetória estética e política de J. Cunha e
sua participação na construção da
identidade visual dos blocos afro e do
carnaval de Salvador
\cite{barata2024jcunha}.

Os dois casos mostram que tecido, figurino,
grafismo, bloco afro e carnaval não
são margens da história da arte: são campos
fundamentais de produção visual
negra em Salvador.

\section{Tincoã e Os Tincoãs}
A exposição \emph{Tincoã} colocava em relação
o pássaro tincoã e o grupo vocal
baiano Os Tincoãs. O movimento curatorial
passava do canto do pássaro ao canto
musical, e daí à memória e ancestralidade.

Essa experiência fixou uma regra para estas
notas: em Salvador, arte visual,
música, religião, artesanato e memória
aparecem frequentemente no mesmo campo
cultural; separá-los por mídia pode esconder
mais do que explicar. O próprio
projeto de Mateus Aleluia descreve Os Tincoãs
como pioneiros em levar de forma
consistente à música popular brasileira
universos poéticos ligados ao
Candomblé e à Umbanda
\cite{nacoes-candomble-projeto}.


\section{Orun Aiye e Jorge Washington}
O curta de animação \emph{Orun Aiye---A
criação do mundo}, dirigido por
Jamile Coelho e Cintia Maria, apareceu no
MUNCAB como outra ponte entre arte e
cosmologia. A produção em stop motion
trabalha narrativas de criação ligadas
ao universo iorubá. Nos créditos, Carlinhos
Brown interpreta Oxalá, Jorge
Washington interpreta Orunmilá, João Miguel
interpreta Olodumaré e Carlos
Betão interpreta Bira
\cite{orun-aiye-sepromi}.

Jorge Washington é ator do Bando de Teatro
Olodum, companhia negra criada em
Salvador em 1990
\cite{bando-jorge-washington}. O encontro do
seu nome no filme
é portanto também uma conexão entre teatro
negro baiano, cinema de animação e
repertório religioso afro-brasileiro.

\section{Casa do Benin e fluxo/refluxo}
A Casa do Benin não deve ser lida apenas como
um prédio que contém objetos
africanos. A instituição foi criada em 1988
para materializar intercâmbios
Bahia--Benin, em paralelo com a Casa do
Brasil em Ouidah
\cite{casa-benin-visite}.

Pierre Verger fornece a chave histórica do
fluxo e refluxo entre as duas
margens \cite{verger2021fluxo}; a intervenção
arquitetônica de Lina Bo Bardi
faz com que o próprio edifício participe
dessa reflexão.

\section{O navio da Liberdade}
Emo de Medeiros, em \emph{O navio da
Liberdade} (2025), reelabora o navio
atlântico como imagem de liberdade em vez de
deportação. A obra pertence ao
projeto \emph{Fatumbia}, que o próprio
artista relaciona a Pierre Fatumbi
Verger e à travessia entre Bahia e África
\cite{emomedeiros-fatumbia}.

A obra é contemporânea e inventada; não deve
ser lida como fotografia
documental de um navio histórico.

\section{Mirongar e os limites da
visibilidade}
O texto \emph{Mirongar}, na Casa do Benin,
propunha trabalhar imagens entre o
que se pode mostrar e aquilo que precisa
permanecer guardado. A formulação
produz uma questão estética e ética
importante: tornar algo visível não é o
mesmo que torná-lo disponível para consumo.

Essa questão ecoa o material religioso
encontrado no MUNCAB e no livro do Opô
Afonjá, onde determinados ritos e
conhecimentos preservam deliberadamente
limites de acesso.

\bibliography{my}
\bibliographystyle{amsalpha}

\end{document}
